\documentclass[11pt,a4paper]{article}

\usepackage[margin=2.2cm]{geometry}
\usepackage[T1]{fontenc}
\usepackage{lmodern}
\usepackage{microtype}
\usepackage{amsmath,amssymb,bm}
\usepackage{booktabs}
\usepackage{array}
\usepackage{graphicx}
\usepackage{caption}
\usepackage{float}
\usepackage{enumitem}
\usepackage{fancyvrb}
\usepackage[hidelinks]{hyperref}
\usepackage{xurl}

\captionsetup{font=small,labelfont=bf}
\setlist{nosep}

% Every number comes from results/numbers.tex, written by `python -m src.run`; a mistyped key
% prints a bold ??.
\makeatletter
\newcommand{\resdef}[2]{\@namedef{res@#1}{#2}}
\newcommand{\res}[1]{\@ifundefined{res@#1}{\textbf{??}}{\@nameuse{res@#1}}}
\newcommand{\resrows}[1]{\@ifundefined{res@#1}{\textbf{??} \\}{\@nameuse{res@#1}}}% table rows
\makeatother
% Generated by python -m src.run from results/*.json; do not edit.
\resdef{p1-orth}{\texttt{4.0e-15}}
\resdef{p1-trace}{13.0000}
\resdef{p1-trace-diff}{\texttt{1.8e-15}}
\resdef{p1-lib-eigenvalues}{\texttt{3.6e-15}}
\resdef{p1-lib-eigenvectors}{\texttt{9.1e-15}}
\resdef{p1-loading-rows}{\texttt{alcohol} & $0.1576$ & $0.4941$ \\ \texttt{malic\_acid} & $-0.2524$ & $0.1671$ \\ \texttt{ash} & $-0.0166$ & $0.3358$ \\ \texttt{alcalinity\_of\_ash} & $-0.2351$ & $-0.0230$ \\ \texttt{magnesium} & $0.1429$ & $0.2817$ \\ \texttt{total\_phenols} & $0.3899$ & $0.0716$ \\ \texttt{flavanoids} & $0.4260$ & $0.0057$ \\ \texttt{nonflavanoid\_phenols} & $-0.2868$ & $0.0486$ \\ \texttt{proanthocyanins} & $0.3086$ & $0.0059$ \\ \texttt{color\_intensity} & $-0.0889$ & $0.5263$ \\ \texttt{hue} & $0.3012$ & $-0.2840$ \\ \texttt{od280/od315\_of\_diluted\_wines} & $0.3778$ & $-0.1814$ \\ \texttt{proline} & $0.2884$ & $0.3703$ \\}
\resdef{p1-top-pc1}{\texttt{flavanoids} ($0.4260$), \texttt{total\_phenols} ($0.3899$) and \texttt{od280/od315\_of\_diluted\_wines} ($0.3778$)}
\resdef{p1-top-pc2}{\texttt{color\_intensity} ($0.5263$), \texttt{alcohol} ($0.4941$) and \texttt{proline} ($0.3703$)}
\resdef{n-train}{124}
\resdef{n-test}{54}
\resdef{p2-k-elbow}{3}
\resdef{p2-R-elbow}{0.6595}
\resdef{p2-tail-elbow}{0.3405}
\resdef{p2-k80}{5}
\resdef{p2-k90}{7}
\resdef{p2-k95}{10}
\resdef{p2-kaiser}{3}
\resdef{p2-identity}{\texttt{3.6e-15}}
\resdef{p2-rows}{1 & $0.3573$ & $0.3573$ & $8.3550$ & $8.5455$ \\ 2 & $0.1921$ & $0.5494$ & $5.8578$ & $6.1120$ \\ 3 & $0.1101$ & $0.6595$ & $4.4270$ & $4.5980$ \\ 4 & $0.0725$ & $0.7320$ & $3.4844$ & $3.8979$ \\ 5 & $0.0697$ & $0.8017$ & $2.5779$ & $3.0326$ \\ 6 & $0.0534$ & $0.8551$ & $1.8835$ & $2.4541$ \\ 7 & $0.0456$ & $0.9007$ & $1.2913$ & $1.8367$ \\ 8 & $0.0242$ & $0.9248$ & $0.9773$ & $1.3861$ \\ 9 & $0.0204$ & $0.9452$ & $0.7120$ & $1.0353$ \\ 10 & $0.0198$ & $0.9650$ & $0.4550$ & $0.7532$ \\ 11 & $0.0169$ & $0.9819$ & $0.2359$ & $0.4497$ \\ 12 & $0.0109$ & $0.9927$ & $0.0947$ & $0.1726$ \\ 13 & $0.0073$ & $1.0000$ & $0.0000$ & $0.0000$ \\}
\resdef{p3-rows}{V1 (centred only) & $0.9978$ & $0.0021$ & $0.9998$ & \texttt{proline} & $0.7222$ \\ V2 (min-max) & $0.4271$ & $0.1841$ & $0.6113$ & \texttt{flavanoids} & $0.9815$ \\ V3 (z-score) & $0.3573$ & $0.1921$ & $0.5494$ & \texttt{flavanoids} & $0.9444$ \\}
\resdef{p3-raw-list}{\texttt{proline} ($94906.7109$), \texttt{magnesium} ($232.0715$) and \texttt{alcalinity\_of\_ash} ($11.4713$)}
\resdef{p3-raw-share}{0.9973}
\resdef{p3-V1-pc1}{\texttt{proline}}
\resdef{p3-V1-pc1-loading}{0.9998}
\resdef{p3-V1-pc2}{\texttt{magnesium}}
\resdef{p3-V1-confused}{1 and 2}
\resdef{p3-V1-confused-errors}{11}
\resdef{p3-V1-errors}{15}
\resdef{p3-V2-lead}{2}
\resdef{p3-V1-acc}{0.7222}
\resdef{p3-V2-acc}{0.9815}
\resdef{p3-V3-acc}{0.9444}
\resdef{p3-V1-r1}{0.9978}
\resdef{p3-V1-r12}{0.9998}
\resdef{p3-V3-r1}{0.3573}
\resdef{p3-V3-r12}{0.5494}
\resdef{p4-baseline}{0.9815}
\resdef{p4-baseline-correct}{53}
\resdef{p4-threshold}{0.9630}
\resdef{p4-one-sample}{0.0185}
\resdef{p4-rows}{1 & $0.8704$ & $0.8519$ & $0.8519$ \\ 2 & $0.9444$ & $0.9074$ & $0.8889$ \\ 3 & $0.9815$ & $0.9074$ & $0.8704$ \\ 5 & $1.0000$ & $0.9815$ & $0.9074$ \\ 8 & $0.9815$ & $0.9815$ & $0.9630$ \\}
\resdef{p4-f-rows}{1 & \texttt{flavanoids} & $187.6858$ \\ 2 & \texttt{proline} & $149.4335$ \\ 3 & \texttt{od280/od315\_of\_diluted\_wines} & $126.6191$ \\ 4 & \texttt{alcohol} & $100.3464$ \\ 5 & \texttt{color\_intensity} & $83.9793$ \\ 6 & \texttt{hue} & $72.2673$ \\ 7 & \texttt{total\_phenols} & $59.2558$ \\ 8 & \texttt{alcalinity\_of\_ash} & $26.0528$ \\ 9 & \texttt{malic\_acid} & $22.5828$ \\ 10 & \texttt{proanthocyanins} & $19.9799$ \\ 11 & \texttt{nonflavanoid\_phenols} & $17.9053$ \\ 12 & \texttt{ash} & $11.5760$ \\ 13 & \texttt{magnesium} & $7.1124$ \\}
\resdef{p4-chosen-rows}{3 & \texttt{flavanoids}, \texttt{proline}, \texttt{od280/od315\_of\_diluted\_wines} & \texttt{flavanoids}, \texttt{alcohol}, \texttt{alcalinity\_of\_ash} & 1 \\ \addlinespace 8 & \texttt{flavanoids}, \texttt{proline}, \texttt{od280/od315\_of\_diluted\_wines}, \texttt{alcohol}, \texttt{color\_intensity}, \texttt{hue}, \texttt{total\_phenols}, \texttt{alcalinity\_of\_ash} & \texttt{flavanoids}, \texttt{alcohol}, \texttt{alcalinity\_of\_ash}, \texttt{proanthocyanins}, \texttt{magnesium}, \texttt{ash}, \texttt{hue}, \texttt{proline} & 5 \\}
\resdef{p4-wrapper-k}{3, 5 and 8}
\resdef{p4-wrapper-cv}{0.9517, 0.9597 and 1.0000}
\resdef{p4-wrapper-test}{0.8704, 0.9074 and 0.9630}
\resdef{p4-se-at}{0.9500}
\resdef{p4-se}{0.0297}
\resdef{p4-gap}{6}
\resdef{p4-gap-methods}{PCA and the wrapper}
\resdef{p4-gap-k}{3}
\resdef{p4-within-pca}{$k = 3$}
\resdef{p4-within-filter}{$k = 5$}
\resdef{p4-within-wrapper}{$k = 8$}
\resdef{p4-shared-3}{\texttt{flavanoids}}
\resdef{p4-shared-8}{\texttt{flavanoids}, \texttt{proline}, \texttt{alcohol}, \texttt{hue} and \texttt{alcalinity\_of\_ash}}

\newcommand{\fig}[2][\linewidth]{\includegraphics[width=#1]{../figures/#2.pdf}}

\newcommand{\T}{\top}
\newcommand{\x}{\mathbf{x}}
\newcommand{\z}{\mathbf{z}}
\newcommand{\uu}{\mathbf{u}}
\newcommand{\cc}{\mathbf{c}}
\newcommand{\xbar}{\bar{\mathbf{x}}}
\newcommand{\tr}{\operatorname{tr}}
\newcommand{\norm}[1]{\lVert #1 \rVert}
\newcommand{\kelbow}{k_{\mathrm{elbow}}}
\newcommand{\Jtrain}{J_{\mathrm{train}}}
\newcommand{\Jtest}{J_{\mathrm{test}}}

\title{\vspace{-1.5cm}CS F320 Foundations of Data Science\\[2pt]
  \Large Assignment 2: Principal Component Analysis}
\author{%
  \begin{tabular}{ll}
    Chaitanya Keyal & 2023B4A70727H \\
    Abhishek Verma & 2023B4A71256H \\
    Palash Jain & 2023A2B51271H
  \end{tabular}}
\date{November 2026}

\begin{document}
\maketitle

\section*{Implementation}

Everything is written in Python~3.12. The covariance matrix, its eigendecomposition, the
projections and reconstructions, the three preprocessing variants, the F-scores and sequential
forward selection are our own NumPy code. scikit-learn is used only where the assignment names
it: \texttt{PCA} for the verification in P1, \texttt{KNeighborsClassifier} in P3, and
\texttt{LogisticRegression} with \texttt{StratifiedKFold} cross-validation in P4. The code, the
data and every output are at
\begin{center}\url{https://github.com/Chaitanya-Keyal/principal-component-analysis}\end{center}
\texttt{src/p1.py} to \texttt{src/p4.py} hold one problem each, \texttt{src/common/pca.py} the
shared PCA code, and \texttt{python -m src.run} regenerates every number and figure in this
report from the two provided CSV files, which are loaded as given.

The conventions C1 to C6 hold throughout. Every statistic is computed on the
$N = \res{p2-n-train}$ training samples and applied unchanged to the $\res{p2-n-test}$ test
samples. Standard deviations and $S$ use the divisor $N - 1$. Eigenvalues and eigenvectors come
from \texttt{numpy.linalg.eigh}, sorted in descending order, and every eigenvector, including
those returned by scikit-learn, is flipped so that its entry of largest absolute value is
positive. Numbers are given to 4 decimal places, and the residuals that are compared with
$10^{-8}$ in scientific notation. Nothing is random except the cross-validation folds of P4,
which use the seed the assignment fixes (\texttt{random\_state=0}).

%==============================================================================================
\section{P1: PCA from first principles}

\paragraph{Derivation.} Let $\uu$ be a unit vector. The projections $\uu^\T\x_n$ of the data
have mean $\uu^\T\xbar$, and because $\uu^\T\mathbf{a} = \mathbf{a}^\T\uu$ for any vector
$\mathbf a$, each squared deviation is
$(\uu^\T\x_n - \uu^\T\xbar)^2 = \uu^\T(\x_n - \xbar)(\x_n - \xbar)^\T\uu$. Summing over $n$ and
dividing by $N - 1$, the variance of the projections is
\[
  \frac{1}{N-1}\sum_{n=1}^{N}\bigl(\uu^\T\x_n - \uu^\T\xbar\bigr)^2 = \uu^\T S\,\uu,
  \qquad
  S = \frac{1}{N-1}\sum_{n=1}^{N}(\x_n - \xbar)(\x_n - \xbar)^\T .
\]
The first principal direction maximises $\uu^\T S\uu$ subject to $\uu^\T\uu = 1$. With the
Lagrangian $\mathcal L(\uu, \lambda) = \uu^\T S\uu + \lambda(1 - \uu^\T\uu)$,
\[
  \nabla_{\uu}\mathcal L = 2S\uu - 2\lambda\uu = \mathbf 0
  \quad\Longrightarrow\quad S\uu = \lambda\uu ,
\]
so every stationary point is an eigenvector of $S$. Multiplying on the left by $\uu^\T$ gives
$\uu^\T S\uu = \lambda$: the variance along an eigenvector is its eigenvalue, and the maximum is
the largest eigenvalue $\lambda_1$, reached at its eigenvector $\uu_1$.

For the second direction we add the constraint $\uu^\T\uu_1 = 0$ with a multiplier $\eta$, which
gives $2S\uu - 2\lambda\uu + \eta\,\uu_1 = \mathbf 0$. Multiplying on the left by $\uu_1^\T$, the
first term is $2\uu_1^\T S\uu = 2(S\uu_1)^\T\uu = 2\lambda_1\uu_1^\T\uu = 0$ because $S$ is
symmetric, the second is $0$ by the constraint, and the third is $\eta$, so $\eta = 0$ and again
$S\uu = \lambda\uu$. Among the eigenvectors orthogonal to $\uu_1$ the largest variance is
$\lambda_2$, at $\uu_2$, and the same argument with one multiplier per earlier component gives
$\uu_3, \dots, \uu_D$.

Finally, with $S = U\Lambda U^\T$ and $U^\T U = I$, $\tr(S) = \tr(\Lambda U^\T U) =
\sum_j\lambda_j$. The trace is also the sum of the per-feature variances $\sum_j s_j^2$, and
after standardisation every $s_j^2 = 1$, so $\sum_j\lambda_j = D = 13$.

\paragraph{Implementation.} The 13 training features are standardised as $z_{nj} =
(x_{nj} - \mu_j)/s_j$. \texttt{covariance} centres the data and forms
$S = \tilde X^\T\tilde X/(N - 1)$, and \texttt{eigendecomposition} calls
\texttt{numpy.linalg.eigh}, reverses its ascending order and applies the sign rule. The class
column is not read by this problem. The program prints the eigenvalues as follows.

\begin{center}
\begin{minipage}{0.32\linewidth}
\VerbatimInput[fontsize=\small]{../results/p1.txt}
\end{minipage}
\end{center}

\paragraph{Checks.} Table~\ref{tab:p1-checks} gives the two checks on our own decomposition and
the comparison with \texttt{PCA(n\_components=13, svd\_solver="full")} fitted on the same
standardised data, after the sign rule has been applied to the rows of \texttt{components\_}.
$\tr(S) = \res{p1-trace}$, as derived above. All four residuals are far below $10^{-8}$, so $U$
is orthonormal, the eigenvalues sum to the trace, and our eigenvalues and eigenvectors agree
with the library's.

\begin{table}[H]
  \centering\small
  \begin{tabular}{lc}
    \toprule
    Quantity & Value \\
    \midrule
    $\norm{U^\T U - I}_F$ & \res{p1-orth} \\
    $\tr(S)$ & $\res{p1-trace}$ \\
    $\bigl|\tr(S) - \sum_j\lambda_j\bigr|$ & \res{p1-trace-diff} \\
    $\max_j\bigl|\lambda_j - \texttt{explained\_variance\_}_j\bigr|$ & \res{p1-lib-eigenvalues} \\
    $\max_{ij}\bigl|U_{ij} - (\texttt{components\_}^\T)_{ij}\bigr|$ & \res{p1-lib-eigenvectors} \\
    \bottomrule
  \end{tabular}
  \caption{Checks on the decomposition and its largest absolute differences from scikit-learn.}
  \label{tab:p1-checks}
\end{table}

\paragraph{Loadings.} Table~\ref{tab:p1-loadings} lists $\uu_1$ and $\uu_2$. The three features
with the largest $|\text{loading}|$ on PC1 are \res{p1-top-pc1}, and on PC2 they are
\res{p1-top-pc2}.

\begin{table}[H]
  \centering\small
  \begin{tabular}{lrr}
    \toprule
    Feature & $\uu_1$ (PC1) & $\uu_2$ (PC2) \\
    \midrule
    \resrows{p1-loading-rows}
    \bottomrule
  \end{tabular}
  \caption{Loadings of the first two principal components.}
  \label{tab:p1-loadings}
\end{table}

%==============================================================================================
\section{P2: Variance captured vs.\ number of components}

\paragraph{Method.} With $r_j = \lambda_j/\sum_i\lambda_i$ and $R(k) = \sum_{j\le k} r_j$, we
report $k_\tau = \min\{k : R(k) \ge \tau\}$ for $\tau \in \{0.80, 0.90, 0.95\}$, the Kaiser count
(components with $\lambda_j > \tr(S)/D$) and $\kelbow = \max\{k : r_k \ge 0.10\}$. A standardised
point is projected and reconstructed as $\hat\z_n = U_kU_k^\T\z_n$ with $U_k = [\uu_1, \dots,
\uu_k]$, and the reconstruction error of a set of $n_{\mathcal D}$ points is
$J_{\mathcal D}(k) = \frac{1}{n_{\mathcal D} - 1}\sum_n\norm{\z_n - \hat\z_n}^2$. Test points are
standardised with the training mean and standard deviation and use the training $U$.

The eigenvectors form an orthonormal basis, so $\z_n = \sum_{j=1}^{D}(\uu_j^\T\z_n)\uu_j$ and the
rank-$k$ reconstruction keeps the first $k$ terms. The residual is $\z_n - \hat\z_n =
\sum_{j>k}(\uu_j^\T\z_n)\uu_j$, and by orthonormality $\norm{\z_n - \hat\z_n}^2 =
\sum_{j>k}(\uu_j^\T\z_n)^2$. The standardised training data has zero mean, so
$\frac{1}{N-1}\sum_n(\uu_j^\T\z_n)^2 = \uu_j^\T S\uu_j = \lambda_j$, and summing over the training
set
\[
  \Jtrain(k) = \sum_{j=k+1}^{D}\lambda_j, \qquad \frac{\Jtrain(k)}{\tr(S)} = 1 - R(k):
\]
the variance the first $k$ components do not capture is exactly the reconstruction error. For
standardised data $\tr(S)/D = 1$, so the Kaiser criterion keeps the components with
$\lambda_j > 1$.

\paragraph{Results.} Figure~\ref{fig:p2-variance} shows the scree plot and the cumulative
variance. The elbow rule gives $\kelbow = \res{p2-k-elbow}$. The thresholds are reached at
$k_{0.80} = \res{p2-k80}$, $k_{0.90} = \res{p2-k90}$ and $k_{0.95} = \res{p2-k95}$, and the Kaiser
criterion retains $\res{p2-kaiser}$ components. Table~\ref{tab:p2} lists $r_k$, $R(k)$ and both
reconstruction errors, each computed directly from the reconstructions $\hat\z_n$, and
Figure~\ref{fig:p2-reconstruction} plots the two errors. For the identity derived above,
\[
  \max_{k=1,\dots,13}\Bigl|\Jtrain(k) - \sum_{j>k}\lambda_j\Bigr| = \text{\res{p2-identity}},
\]
which is below $10^{-8}$, so the two sides agree.

\begin{figure}[H]
  \centering\fig{p2_variance}
  \caption{Left: explained variance ratio $r_j$ of each component, with the elbow
    $\kelbow = \res{p2-k-elbow}$. Right: cumulative ratio $R(k)$ with the reference levels 0.80,
    0.90 and 0.95.}
  \label{fig:p2-variance}
\end{figure}

\begin{table}[H]
  \centering\small
  \begin{tabular}{rcccc}
    \toprule
    $k$ & $r_k$ & $R(k)$ & $\Jtrain(k)$ & $\Jtest(k)$ \\
    \midrule
    \resrows{p2-rows}
    \bottomrule
  \end{tabular}
  \caption{Explained variance and reconstruction error for every $k$.}
  \label{tab:p2}
\end{table}

\begin{figure}[H]
  \centering\fig[0.7\linewidth]{p2_reconstruction}
  \caption{Reconstruction error of the training and the test set against $k$.}
  \label{fig:p2-reconstruction}
\end{figure}

\paragraph{Discussion.} (i) The elbow rule gives $\kelbow = \res{p2-k-elbow}$ and the Kaiser
criterion also keeps $\res{p2-kaiser}$ components, fewer than $k_{0.80} = \res{p2-k80}$,
$k_{0.90} = \res{p2-k90}$ and $k_{0.95} = \res{p2-k95}$. Since $R(\kelbow) = \res{p2-R-elbow}$,
the elbow keeps every component that explains at least a tenth of the variance but leaves
$\res{p2-tail-elbow}$ of the total variance in the discarded tail, so it suits a compact
description of the data and not an accurate reconstruction. (ii) $\Jtest(k)$ is non-increasing in
$k$: the residual expression holds for any point, so $\Jtest(k) - \Jtest(k+1) =
\frac{1}{n_{\mathcal D} - 1}\sum_n(\uu_{k+1}^\T\z_n)^2 \ge 0$. Figure~\ref{fig:p2-reconstruction}
shows this, with both curves falling at every step and reaching zero at $k = 13$, where
$U_{13}U_{13}^\T = I$ reconstructs every point exactly. The test curve lies above the training
curve for every $k < 13$, because the directions were fitted to the training points.

%==============================================================================================
\section{P3: Effect of preprocessing, and the two-dimensional projection}

\paragraph{Method.} Three variants are fitted on the training set and applied unchanged to the
test set: V1 subtracts the training mean; V2 maps each feature to $[0, 1]$ with the training
minimum and maximum and then subtracts the training mean of the scaled feature; V3 is the
z-score. For each one we compute $S$ and its eigendecomposition with the code of P1, project
onto the first two components, $\cc_n = U_2^\T\x'_n$, and train a 5-nearest-neighbour classifier
(\texttt{KNeighborsClassifier(n\_neighbors=5)}) on the training projections, scored on the test
projections.

Two facts explain the results. If one feature $j$ is multiplied by $a > 1$, then $S_{jj}$ becomes
$a^2S_{jj}$ while the other diagonal entries stay the same, and because $\lambda_1 \ge
\mathbf e_j^\T S\,\mathbf e_j = S_{jj}$,
\[
  r_1 \;\ge\; \frac{S_{jj}}{\tr(S)} = \frac{a^2S_{jj}}{a^2S_{jj} + \sum_{i\ne j}S_{ii}}
  \;\longrightarrow\; 1 \quad (a \to \infty),
\]
so the leading component turns into the feature with the largest scale. For z-scored data,
$S_{ji} = \frac{1}{N-1}\sum_n z_{nj}z_{ni} = \mathrm{Cov}(x_j, x_i)/(s_js_i) = \rho_{ji}$: $S$ is
the correlation matrix, which does not depend on the units of measurement.

\paragraph{Results.} Table~\ref{tab:p3} compares the variants and Figure~\ref{fig:p3} shows the
three projections. For V1, the three features with the largest raw sample variance are
\res{p3-raw-list}.

\begin{table}[H]
  \centering\small
  \begin{tabular}{lccclc}
    \toprule
    Variant & $r_1$ & $r_2$ & $r_1 + r_2$ & largest $|\text{loading}|$ on PC1 & 5-NN test accuracy \\
    \midrule
    \resrows{p3-rows}
    \bottomrule
  \end{tabular}
  \caption{Variance captured by the first two components, the feature that dominates PC1, and
    the 5-NN test accuracy on the two-dimensional projection.}
  \label{tab:p3}
\end{table}

\begin{figure}[H]
  \centering\fig{p3_projections}
  \caption{PC1--PC2 projections under the three preprocessing variants. Filled circles are
    training points and crosses are test points, coloured by class.}
  \label{fig:p3}
\end{figure}

\paragraph{Discussion.} (i) In V1, $\tr(S)$ is the sum of the raw variances and \res{p3-V1-pc1}
alone holds $\res{p3-raw-share}$ of it, so by $r_1 \ge S_{jj}/\tr(S)$ the first component is
essentially the \res{p3-V1-pc1} axis and $r_1 = \res{p3-V1-r1}$. V3 gives every feature variance
1, so $S$ is the correlation matrix and $r_1 = \res{p3-V3-r1}$ measures variation shared by
correlated features, not the units of one of them. (ii) V1 captures $\res{p3-V1-r12}$ of the
variance in two components against $\res{p3-V3-r12}$ for V3, yet it is the worst for the
classification check, with a 5-NN accuracy of $\res{p3-V1-acc}$ against $\res{p3-V3-acc}$: in
Figure~\ref{fig:p3} its plane is spanned almost entirely by \res{p3-V1-pc1} and
\res{p3-V1-pc2}, on which classes 1 and 2 overlap. Variance captured is therefore not a measure
of useful structure, because PCA ranks directions by spread alone, a spread that depends on the
units, while the class information can sit in features with a small raw variance. (iii) We would
use the z-score (V3) before PCA on this data: the thirteen features have unrelated units and
ranges, and standardising makes the components independent of those units. V2 also removes the
unit problem and scores $\res{p3-V2-acc}$ here, but it scales each feature by its two most
extreme training values, so it is sensitive to outliers, and its lead over V3 is only
\res{p3-V2-lead} of the \res{p3-n-test} test samples.

%==============================================================================================
\section{P4: PCA versus feature selection for a downstream classifier}

\paragraph{Method.} The classifier is the same for every feature set, a multinomial logistic
regression, \texttt{LogisticRegression(C=1.0, max\_iter=1000)}, trained on the standardised
training set and scored on the standardised test set. Cross-validation uses the folds of
\texttt{StratifiedKFold(n\_splits=5, shuffle=True, random\_state=0)} on the training set, built
once and reused. For $k \in \{1, 2, 3, 5, 8\}$ the three feature sets are
\begin{itemize}
  \item[(A)] \textbf{PCA}: the first $k$ principal component scores $U_k^\T\z_n$, with $U$ from
        the training set (P1);
  \item[(B)] \textbf{Filter}: the $k$ original features with the largest F-score
        \[
          F_j = \frac{\sum_{c=1}^{C}N_c\,(m_{jc} - m_j)^2\big/(C - 1)}
                     {\sum_{c=1}^{C}\sum_{n\in c}(z_{nj} - m_{jc})^2\big/(N - C)},
        \]
        the between-class variance of feature $j$ over its within-class variance, computed on
        the training set;
  \item[(C)] \textbf{Wrapper}: sequential forward selection. Starting from the empty set, each
        step adds the feature whose addition gives the highest mean 5-fold cross-validated
        accuracy; candidates are tried in increasing column index and the first maximum is
        kept.
\end{itemize}
The baseline is the same classifier on all 13 standardised features. Forward selection is our
own implementation; it returns the same subsets as scikit-learn's
\texttt{SequentialFeatureSelector} for every $k$ in the table (checked in
\texttt{tests/test\_p4.py}). At a step where $k$ features are held it scores $D - k$ candidates,
so running to completion would score $\sum_{k=0}^{D-1}(D - k) = D(D+1)/2 = 91$ subsets, against
$2^{D} - 1 = 8191$ for an exhaustive search; stopping at $k = 8$ scores $13 + 12 + \dots + 6 = 76$
subsets, each with 5 model fits.

\paragraph{Results.} Table~\ref{tab:p4} and Figure~\ref{fig:p4} give the test accuracy of the
three methods. The baseline test accuracy is $\res{p4-baseline}$ ($\res{p4-baseline-correct}$ of
the $\res{p4-n-test}$ test samples).

\begin{table}[H]
  \centering\small
  \begin{tabular}{rccc}
    \toprule
    $k$ & (A) PCA & (B) Filter & (C) Wrapper \\
    \midrule
    \resrows{p4-rows}
    \bottomrule
  \end{tabular}
  \caption{Test accuracy of the classifier on $k$ features from each method.}
  \label{tab:p4}
\end{table}

\begin{figure}[H]
  \centering\fig[0.72\linewidth]{p4_accuracy}
  \caption{Test accuracy against $k$ for the three methods, with the baseline on all 13
    features.}
  \label{fig:p4}
\end{figure}

Table~\ref{tab:p4-f} ranks all 13 features by F-score, and Table~\ref{tab:p4-chosen} lists the
features each selection method chooses, (B) in rank order and (C) in the order they were added.
At $k = 3$ the two methods share \res{p4-shared-3}; at $k = 8$ they share \res{p4-shared-8}.

\begin{table}[H]
  \centering\small
  \begin{tabular}{rlr}
    \toprule
    Rank & Feature & $F_j$ \\
    \midrule
    \resrows{p4-f-rows}
    \bottomrule
  \end{tabular}
  \caption{F-score ranking of the 13 features, best first.}
  \label{tab:p4-f}
\end{table}

\begin{table}[H]
  \centering\small
  \begin{tabular}{r>{\raggedright\arraybackslash}p{0.38\linewidth}%
      >{\raggedright\arraybackslash}p{0.38\linewidth}c}
    \toprule
    $k$ & (B) Filter & (C) Wrapper & Shared \\
    \midrule
    \resrows{p4-chosen-rows}
    \bottomrule
  \end{tabular}
  \caption{Features chosen by the filter and by the wrapper, and how many they share.}
  \label{tab:p4-chosen}
\end{table}

\paragraph{Discussion.} (i) Within one test sample of the baseline, at an accuracy of at least
$\res{p4-threshold}$, PCA arrives at \res{p4-within-pca}, the filter at \res{p4-within-filter}
and the wrapper at \res{p4-within-wrapper}. (ii) The labels give (B) and (C) no visible advantage
in Figure~\ref{fig:p4}: the PCA curve is on or above both at every $k$, and even at $k = 1$ one
component found without labels ($\res{p4-first-pc-acc}$) does at least as well as
\res{p4-first-feature}, the single feature both supervised methods pick
($\res{p4-first-feature-acc}$). That is consistent with how the features are built: a principal
component is a weighted sum of all thirteen features, so $k$ components draw on every feature
while $k$ selected features draw on only $k$, and on this data the directions of largest variance
also separate the classes (P3). The one trace of the labels is in the wrapper, whose
cross-validated accuracy on the training set reaches $\res{p4-sfs-cv-last}$ at $k = 8$ while its
test accuracy is the lowest of the three for every $k \ge 2$, a sign that it tunes the subset to
the $\res{p2-n-train}$ training samples. (iii) With $\res{p4-n-test}$ test samples one sample is
worth $\res{p4-one-sample}$ of accuracy and the binomial standard error of an accuracy near 0.95
is about 0.03, so differences of one or two samples, such as those between PCA and the filter
from $k = 5$ on, cannot rank the methods. Only the larger gaps at $k = 3$ (\res{p4-gap-filter-3}
and \res{p4-gap-wrapper-3} samples between PCA and the two selection methods) point to a real
difference, and even they come from a single train/test split.

\end{document}
